%%%%%%%%%%%%%%%%%%%%%%%%%%%%%%%%%%%%%%%%%%%%%%%%%%%%%
%			HLAVIČKA								%
%%%%%%%%%%%%%%%%%%%%%%%%%%%%%%%%%%%%%%%%%%%%%%%%%%%%%
\documentclass[openany,12pt]{memoir}

\usepackage[utf8]{inputenc} 
\usepackage[czech]{babel}
\usepackage[T1]{fontenc}
\usepackage[top=1.5cm, bottom=2cm, left=2cm, right=2cm]{geometry}  % --> NASTAVENÍ OKRAJŮ
\usepackage{fancyhdr}
\usepackage{graphicx}
%\usepackage{lmodern}
\usepackage{tgschola}  % --- Silnější font
\usepackage{xwatermark}
\usepackage{xcolor}
\usepackage{changepage}
\usepackage{newunicodechar}



%%%%%% Package na zpěvník
\usepackage[full]{leadsheets}%http://mirrors.nic.cz/tex-archive/macros/latex/contrib/leadsheets/leadsheets_en.pdf   --> dokumentace	
\definesongtitletemplate{empty}{} 
\setchords{
format = \bfseries,   %tučné akordy
minor = {mi},% 
input-notation = {german},%
output-notation = {german}%
}
\definesongtitletemplate{empty}{} 

\newlength{\drop}
\newwatermark[allpages,color=red!50,angle=0,scale=2, xpos=0,ypos=0]{\includegraphics[width=5cm]{obr/pozadi2.jpg}} %--> dvojka na pozadí


%%%% Vlastní příkazy
\newcounter{Slokočet}   %Automatické číslování slok
\newcommand{\mezera}{\vspace*{0.5cm}}   %Horizontální odsazení slok
\newcommand{\stred}{5.2cm}   %%% Na zarovnání slok doprostřed, pozn. automatičtější zarovnávání na střed nejde
\newcommand{\refren}{\mezera \noindent \textbf{R:} } %refrén
\newcommand{\sloka}{\mezera \noindent \addtocounter{Slokočet}{1} \arabic{Slokočet}. } 	%sloka, která se automaticky čísluje
\newcommand{\ssloka}{\mezera \noindent}  % vlastní číslo sloky
\newcommand{\carka}{,\:}
\newcommand{\predehra}{\mezera \hspace*{-1.85cm} \textbf{Předehra:} }
\newcommand{\m}[1]{\color{white}{#1}}  %Pro akordy
\newcommand{\ap}{'} %Pro apostrof
\newcommand{\elipsa}{\kern\fontdimen3\font} %Příkaz pro lepší zacházení s výpustkami (=...); je to vpodstatě jen mezera mezi tečkama výpustky

%%%%%%%%%%%%%%%%%%%%%%%%%%%%%%%%%%%%%%%%%%%%%%%%%%%%%%
%			NÁVOD									 %
%%%%%%%%%%%%%%%%%%%%%%%%%%%%%%%%%%%%%%%%%%%%%%%%%%%%%%
%1. Věci v hlavičce IGNOROVAT
%2. Píseň psát do prostoru mezi \begin{song} a \end{song}
%3. další řádek se značí dvěma odsazeníma (= dvakrát stisknout enter)
%4. \refren vždy na začátku refrenu a \sloka na začátek sloky (automaticky se čísluje)
%  \ = alt gr + q ; [/] = alt gr f/g ; {/} = alt gr + b/n; ^ = alt gr + 3 + mezera
%Cokoliv napíšete do ^{  } se bude brát jako akord
%když se toto bude dotýkat nějakého slova (nebude mezi tím a slovem mezera)
%tak se akordy zjeví nad slovem, ale pište to před slovo
%Když se to nedotýká slova, tak akord lítá ve vzduchu a vytiskne se větší mezera
%První možnost je asi preferovanější
%5. Akordy stačí psát jen do první sloky, když se nezmění -- kytaristi to zvládnou



%\usepackage{subfiles}
%</++>

\begin{document}

\pagestyle{empty}
%\addtocontents{toc}{\protect\thispagestyle{empty}}
%\tableofcontents \thispagestyle{empty}\newpage
\newgeometry{top=1.5cm, bottom = 0cm, left = 2cm, right = 2cm}
	
\newpage

%Pro kompilaci je důležitý, aby ten \input a \end{document} zůstali na posledních
%dvou řádcích!!!!


\begin{song}{title=\centering Lišaji \\\normalsize Zuzana Navarová \vspace*{-0.3cm}}  %% sem se napíše jméno songu a autor
\moveright 2.8cm \vbox{      %Varianta č. 1  ---> Jeden sloupec zarovnaný na střed

\sloka
^{G{\color{white}\_\_\_\_\_}}Kroužili ^{D{\color{white}\_\_\_}}lišaji ^{Emi}při samém ^{Hmi7{\color{white}\_}}okraji, ^{Ami}při samém okraji ^{D{\color{white}\_\_\_}}života. 

^{G}Při samém ^{D{\color{white}\_\_\_}}okraji ^{Emi{\color{white}\_}}stíny dva ^{Hmi7{\color{white}\_}}lišají. ^{Ami{\color{white}\_}}Srdce jim tluče a ^{D{\color{white}\_\_\_\_}}třepotá. 
 
\refren
^{Hmi\,\,}Déšť ^{H7{\color{white}\_}}bije do oken, tma ^{Emi{\color{white}\_}}něžně skřípe do skel 

^{Ami}jak v spánku ^{C{\color{white}\_\_\_\_\_}}hlubokém, kdy ^{Ami{\color{white}\_}}kouzla tajou ^{D{\color{white}\_\_\_\_}}voskem. 
 
\sloka
^{G{\color{white}\_\_\_\_\_}}Kroužili ^{D{\color{white}\_\_\_}}lišaji ^{Emi}při samém ^{Hmi7{\color{white}\_}}okraji ^{Ami}po kraji ^{D{\color{white}\_\_}}bota ^{{\color{white}\_\_\_}G}nebota. 

^{Ami\,\,}Hou, ^{Hmi\,\,}hou, ^{Ami}hou a ^{G{\color{white}\_}}hou s ^{Emi{\color{white}\_\_}}písničkou ^{C{\color{white}\_\_\_\_}D\,\,}lišajovou. 
 
 
\phantom{.}
 
\textbf{G\,\,D\,\,Emi\,\,Hmi7\,\,Ami\,\,D\,\,G\,\,D\,\,Emi\,\,Hmi7\,\,Ami\,\,D\,\,G}
 
 \phantom{.}
 
 
\sloka
Kroužili lišaji při samém okraji, při samém okraji života. 

Při samém okraji stíny dva lišají. Srdce jim tluče a třepotá. 
 

\refren
 

\sloka
^{G\,\,}Kéž zavedou ^{D}do ráje ^{Emi{\color{white}\_}}lišaje, ^{Hmi7{\color{white}\_}}lišaje, ^{Ami{\color{white}\_\_}}bude-liž \uv{ó  jé} ^{{\color{white}\_\_\_}G}sobota. 

^{Ami\,\,}Hou, ^{Hmi\,\,}hou, ^{Ami}hou a ^{G}hou s ^{Emi{\color{white}\_\_\_}}písničkou ^{C{\color{white}\_\_\_}D}lišajovou.
 
 \phantom{.}

^{G\,\,}Kéž zavedou ^{D}do ráje ^{Emi{\color{white}\_}}lišaje, ^{Hmi7{\color{white}\_}}lišaje, ^{Ami{\color{white}\_\_\_}}bude-liž \uv{ó  jé} ^{{\color{white}\_\_\_}G}sobota. 

}
\setcounter{Slokočet}{0}
\end{song}



\end{document}
